%% Kernel fixture: a tensor standing at a point of a curve (the threaded loop
%% of Section III, lines 305-327: a closed curve carrying four tensors, one at
%% each of four of its points).
%% Expected records: four ket atoms on row 1 with physical legs up, four
%% invisible junction atoms at the loop's four corners, a closed chain of four
%% named strings through those junctions, and four box atoms whose places are
%% fractions along the chain segment for their own side of the loop.
%% Expected ink: the loop is a closed quadrilateral of four chained straight
%% segments (no via-driven smooth Hobby cycle any more); each box stands ON
%% its segment at the fraction its address names, so the segment enters one
%% side of the box and leaves the other unchanged.  A box is not a junction:
%% nothing about the loop's route or the tensors below it is decided by where
%% the boxes stand.  The four original arclength fractions (0.05, 0.3, 0.55,
%% 0.8) were spaced a uniform quarter-turn apart, so each box keeps the same
%% 0.2 local fraction into its own edge of the chain.
\documentclass{standalone}
\usepackage{tenkz}
\makeatletter
\ExplSyntaxOn
\file_input:n { tenkz-kernel.code.tex }
\ExplSyntaxOff
\makeatother
\begin{document}
\tenkzkernel
\begin{tenkz}[rows={ket}, cols=4, physical=up]
  \tn[at=1 n of (1,1), skin=none, name=jL1]{}
  \tn[at=1 n of (1,4), skin=none, name=jL2]{}
  \tn[at=1 s of (1,4), skin=none, name=jL3]{}
  \tn[at=1 s of (1,1), skin=none, name=jL4]{}
  \tnwire[kind=string, species=g, name=L,
          cross={over at crossing of L and Lb}]{jL1}{jL2}
  \tnwire[kind=string, species=g, name=Lb,
          cross={over at crossing of Lb and Lc}]{jL2}{jL3}
  \tnwire[kind=string, species=g, name=Lc,
          cross={over at crossing of Lc and Ld}]{jL3}{jL4}
  \tnwire[kind=string, species=g, name=Ld,
          cross={over at crossing of Ld and L}]{jL4}{jL1}
  \tn[skin=box, at=on L 0.2]{$1$}
  \tn[skin=box, at=on Lb 0.2]{$2$}
  \tn[skin=box, at=on Lc 0.2]{$3$}
  \tn[skin=box, at=on Ld 0.2]{$4$}
\end{tenkz}
% A closed string still carries: the wound flux ring is closed by its
% homotopy class, and the box stands on it at the fraction its address
% names -- the pin READING-strings.md keeps on this fixture.
\begin{tenkz}[lattice={2x2}, surface=torus]
  \tnwire[kind=string, species=flux, name=L, closed, wind={1,0}]
  \tn[skin=box, at=on L 0.5]{$5$}
\end{tenkz}
\end{document}
